% day2_seg3_kalstats/proc_noise.tex
\begin{frame}
  \onslide<1->{
	\frametitle{Process Noise} 
	\framesubtitle{Characteristics}
}
\begin{itemize}
  \onslide<2->{
 \item The process noise $w(k)$ is assumed to be a white, zero-mean, Gaussian sequence.
}
  \onslide<3->{
 \item Although the continuous-time form of white noise does not really exist, it is sufficient in the 
 discrete-time case  to consider a random process sequence $w(k)$ that is uncorrelated from 
 step to step: $\expect{w(j)w\trans (k)}=0$ 
 if $j\neq k$
}
  \onslide<4->{
 \item This can be achieved if $w(k)$ is the result of sampling a Gauss-Markov process whose exponential 
autocorrelation function rapidly approaches 0 before the next sample is taken 
}
  \onslide<5->{
\item  Process noise is is thought to arise from unmodelled or negligible states that in aggregate 
behave is a normally-distributed fashion 
}
\end{itemize}
\end{frame}


\begin{frame}
  \onslide<1->{
	\frametitle{Process Noise} 
	\framesubtitle{Process Noise Covariance Matrix $Q$}
}
\begin{itemize}
  \onslide<2->{
\item Process noise covariance matrix $Q$ affects the Kalman gain indirectly through the 
covariance propagation equation for $\covmat{k}{-}$:
\[\altcovextrapeq\]
}
  \onslide<3->{
\item This updated error covariance is then included in both the "numerator" and "denominator" of the Kalman gain:
\[\kalgaineq\]
}
  \onslide<4->{
\item If $Q$ is large relative to $R$ then $K\approx I$ causing $\hat{x}^{-}$ to cancel out and 
the updated state estimate $\hat{x}^{+}$ will be nearly identical to the measurement $z$
}
\end{itemize}
\end{frame}

\begin{frame}
  \onslide<1->{
	\frametitle{Process Noise} 
	\framesubtitle{Importance for Filter Tuning}
}
\begin{itemize}
  \onslide<2->{
\item Since $w(k)$ is an unknown, random forcing function, its parameters have to be selected in a 
process called \textbf{filter tuning}  
}
  \onslide<3->{
\item  Filter tuning means selecting appropriate values for the entries of $Q$ 
}
  \onslide<4->{
\item Techniques to augment the state vector are also used to tune the process noise
}
  \onslide<5->{
\item  Process noise covariance $Q$ is one of two key tuning parameters, the other being measurement noise covariance $R$
}
\end{itemize}
\end{frame}

